\clearpage
\section{罚方法}
\label{sec:ch1-penalty-method}

微可压缩流体的定常线性化方程为
\begin{equation}
-\nu\Delta u_\varepsilon
-\frac{1}{\varepsilon}\operatorname{grad}\operatorname{div}u_\varepsilon=f
\quad\text{在 }\Omega\text{ 中},
\label{eq:ch1-penalty-momentum}
\end{equation}
\begin{equation}
u_\varepsilon=0\quad\text{在 }\partial\Omega\text{ 上},
\label{eq:ch1-penalty-boundary}
\end{equation}
其中 $\varepsilon>0$ 是``小''参数. 方程 \eqref{eq:ch1-penalty-momentum}--\eqref{eq:ch1-penalty-boundary} 也是弹性理论中的定常 Lamé 方程. 这些方程还有一种与优化理论和变分法有关的解释: Stokes 问题对应于在满足 $\operatorname{div}v=0$ 的 $v\in\mathbf H_0^1(\Omega)$ 中极小化泛函 $\tfrac12\lVert v\rVert^2-(f,v)$. 若把方程 $\operatorname{div}v=0$ 看作约束, 则 $p$ 是与此约束相联系的 Lagrange 乘子(见 I. Ekeland 和 R. Temam [1]). 从变分法观点看, 自然要引入问题的罚形式: 在 $v\in\mathbf H_0^1(\Omega)$ 中极小化 $\tfrac12\lVert v\rVert^2+(1/2\varepsilon)\lvert\operatorname{div}v\rvert^2-(f,v)$; 该问题的 Euler 方程就是 \eqref{eq:ch1-penalty-momentum}--\eqref{eq:ch1-penalty-boundary}.

R. Temam [2a], [2b] 在完整方程(非线性含时 Navier--Stokes 方程)的一般层次上把罚方法引入 Navier--Stokes 方程. 罚方法在有限元方法和流动问题数值计算中的应用很多; 见 T. Kawai [1] 及其参考文献.

本节将证明, 对固定的 $\varepsilon>0$, 方程 \eqref{eq:ch1-penalty-momentum}--\eqref{eq:ch1-penalty-boundary} 有唯一解 $u_\varepsilon$, 且当 $\varepsilon\to0$ 时 $u_\varepsilon$ 收敛到 Stokes 方程的解 $u$. 第 \ref{subsec:ch1-penalty-convergence} 节说明 $u_\varepsilon$ 与 $u$ 的关系, 第 \ref{subsec:ch1-penalty-asymptotic-expansion} 节给出 $\varepsilon\to0$ 时 $u_\varepsilon$ 的渐近展开, 第 \ref{subsec:ch1-penalty-numerical-algorithms} 节再说明如何利用该渐近展开计算小 $\varepsilon$ 情形下的 $u_\varepsilon$.

\subsection{从 \texorpdfstring{$u_\varepsilon$}{u epsilon} 到 \texorpdfstring{$u$}{u} 的收敛}
\label{subsec:ch1-penalty-convergence}

\MBSourceStatementNumber{theorem}{1}
\begin{Theorem}
\label{thm:ch1-penalty-convergence}
设 $\Omega$ 是 $\mathbb R^n$ 中的有界 Lipschitz 区域. 对固定的 $\varepsilon>0$, 存在唯一的 $u_\varepsilon\in\mathbf H_0^1(\Omega)$ 满足 \eqref{eq:ch1-penalty-momentum}--\eqref{eq:ch1-penalty-boundary}. 当 $\varepsilon\to0$ 时,
\begin{equation}
u_\varepsilon\longrightarrow u
\quad\text{在 }\mathbf H_0^1(\Omega)\text{ 的范数中},
\label{eq:ch1-penalty-velocity-convergence}
\end{equation}
\begin{equation}
-\frac{\operatorname{div}u_\varepsilon}{\varepsilon}
\longrightarrow p
\quad\text{在 }L^2(\Omega)\text{ 的范数中},
\label{eq:ch1-penalty-pressure-convergence}
\end{equation}
其中 $u$ 和 $p$ 由 \eqref{eq:ch1-stokes-variational-problem}--\eqref{eq:ch1-stokes-weak-boundary} 定义, 并且
\begin{equation}
\int_\Omega p(x)\,dx=0.
\label{eq:ch1-penalty-pressure-zero-mean}
\end{equation}
\end{Theorem}

\begin{Proof}
容易证明问题 \eqref{eq:ch1-penalty-momentum}--\eqref{eq:ch1-penalty-boundary} 等价于下述变分问题: 求 $u_\varepsilon\in\mathbf H_0^1(\Omega)$, 使
\begin{equation}
\nu((u_\varepsilon,v))
+\frac{1}{\varepsilon}(\operatorname{div}u_\varepsilon,\operatorname{div}v)
=(f,v),
\qquad\forall v\in\mathbf H_0^1(\Omega).
\label{eq:ch1-penalty-variational-problem}
\end{equation}
事实上, 若 $u_\varepsilon$ 满足 \eqref{eq:ch1-penalty-momentum}--\eqref{eq:ch1-penalty-boundary}, 则 $u_\varepsilon\in\mathbf H_0^1(\Omega)$ 且对每个 $v\in\mathcal D(\Omega)$ 满足 \eqref{eq:ch1-penalty-variational-problem}. 反之, 若 $u_\varepsilon\in\mathbf H_0^1(\Omega)$ 是 \eqref{eq:ch1-penalty-variational-problem} 的解, 则 $u_\varepsilon$ 在分布意义下满足 \eqref{eq:ch1-penalty-momentum}, 并在迹定理意义下满足 \eqref{eq:ch1-penalty-boundary}.

满足 \eqref{eq:ch1-penalty-variational-problem} 的 $u_\varepsilon$ 的存在性与唯一性由投影定理得到: 在定理 \ref{thm:ch1-abstract-projection} 中取
\[
W=\mathbf H_0^1(\Omega),
\qquad a(u,v)=\nu((u,v))+\frac{1}{\varepsilon}(\operatorname{div}u,\operatorname{div}v),
\qquad \langle\ell,v\rangle=(f,v).
\]
$a$ 的强制性以及 $a$ 和 $\ell$ 的连续性都是显然的.

为证明 \eqref{eq:ch1-penalty-velocity-convergence}, 从 \eqref{eq:ch1-penalty-momentum} 中减去 \eqref{eq:ch1-stokes-weak-momentum}, 得
\begin{equation}
-\nu\Delta(u_\varepsilon-u)
-\frac{1}{\varepsilon}\operatorname{grad}\operatorname{div}u_\varepsilon
=\operatorname{grad}p,
\label{eq:ch1-penalty-error-equation}
\end{equation}
因而
\begin{equation}
\nu((u_\varepsilon-u,v))
+\frac{1}{\varepsilon}(\operatorname{div}u_\varepsilon,\operatorname{div}v)
=-(p,\operatorname{div}v),
\qquad\forall v\in\mathbf H_0^1(\Omega).
\label{eq:ch1-penalty-error-variational}
\end{equation}
对 $v\in\mathcal D(\Omega)$, 方程 \eqref{eq:ch1-penalty-error-variational} 容易从 \eqref{eq:ch1-penalty-error-equation} 得到; 由连续性论证, 它对每个 $v\in\mathbf H_0^1(\Omega)$ 成立.

在 \eqref{eq:ch1-penalty-error-variational} 中取 $v=u_\varepsilon-u$, 得
\[
\nu\lVert u_\varepsilon-u\rVert^2
+\frac{1}{\varepsilon}\lvert\operatorname{div}u_\varepsilon\rvert^2
=-(p,\operatorname{div}u_\varepsilon)
\leq\lvert p\rvert\lvert\operatorname{div}u_\varepsilon\rvert
\leq\frac{1}{2\varepsilon}\lvert\operatorname{div}u_\varepsilon\rvert^2
+\frac{\varepsilon}{2}\lvert p\rvert^2,
\]
从而
\begin{equation}
\nu\lVert u_\varepsilon-u\rVert^2
+\frac{1}{2\varepsilon}\lvert\operatorname{div}u_\varepsilon\rvert^2
\leq\frac{\varepsilon}{2}\lvert p\rvert^2.
\label{eq:ch1-penalty-basic-estimate}
\end{equation}
这证明了 \eqref{eq:ch1-penalty-velocity-convergence}. 因而, 由 \eqref{eq:ch1-penalty-error-equation},
\begin{equation}
\frac{\partial}{\partial x_i}
\left(\frac{\operatorname{div}u_\varepsilon}{\varepsilon}\right)
\longrightarrow-\frac{\partial p}{\partial x_i}
\quad\text{在 }H^{-1}(\Omega)\text{ 的范数中},
\quad i=1,\ldots,n,
\label{eq:ch1-penalty-pressure-derivative-convergence}
\end{equation}
因为由 \eqref{eq:ch1-penalty-velocity-convergence}, $\Delta(u_\varepsilon-u)$ 在 $H^{-1}(\Omega)$ 中收敛到 $0$.

根据下面的引理,
\begin{equation}
\left\lvert p+\frac{\operatorname{div}u_\varepsilon}{\varepsilon}\right\rvert
\leq\sum_{i=1}^{n}
\left\lVert\frac{\partial}{\partial x_i}
\left(p+\frac{\operatorname{div}u_\varepsilon}{\varepsilon}\right)
\right\rVert_{H^{-1}(\Omega)},
\label{eq:ch1-penalty-pressure-lemma-application}
\end{equation}
这是因为 $u_\varepsilon$ 在 $\partial\Omega$ 上为零且 \eqref{eq:ch1-penalty-pressure-zero-mean} 成立, 所以
\[
\int_\Omega\left(p+\frac{\operatorname{div}u_\varepsilon}{\varepsilon}\right)dx=0.
\]
于是 \eqref{eq:ch1-penalty-pressure-convergence} 得证.
\end{Proof}

\MBSourceStatementNumber{lemma}{1}
\begin{Lemma}
\label{lem:ch1-negative-norm-pressure-estimate}
设 $\Omega$ 是 $\mathbb R^n$ 中的有界 Lipschitz 区域. 则存在只依赖于 $\Omega$ 的常数 $c=c(\Omega)$, 使对每个 $\sigma\in L^2(\Omega)$,
\begin{equation}
\lvert\sigma\rvert_{L^2(\Omega)}
\leq c(\Omega)\left\{
\left\lvert\int_\Omega\sigma\,dx\right\rvert
+\sum_{i=1}^{n}
\left\lVert\frac{\partial\sigma}{\partial x_i}\right\rVert_{H^{-1}(\Omega)}
\right\}.
\label{eq:ch1-negative-norm-pressure-estimate}
\end{equation}
\end{Lemma}

\begin{Proof}
用 $[\sigma]$ 表示 \eqref{eq:ch1-negative-norm-pressure-estimate} 右端大括号内的量. $[\sigma]$ 是 $L^2(\Omega)$ 上的范数: 它显然是半范数; 若 $[\sigma]=0$, 则由 $\partial\sigma/\partial x_i=0$, $i=1,\ldots,n$, 知 $\sigma$ 是常数, 又由 $\int_\Omega\sigma\,dx=0$ 知该常数为零.

显然存在常数 $c'=c'(\Omega)$, 使
\begin{equation}
[\sigma]\leq c'(\Omega)\lvert\sigma\rvert_{L^2(\Omega)},
\qquad\forall\sigma\in L^2(\Omega).
\label{eq:ch1-negative-norm-upper-bound}
\end{equation}
若证明 $L^2(\Omega)$ 对范数 $[\sigma]$ 完备, 则由闭图定理, $[\sigma]$ 与 $\lvert\sigma\rvert$ 是 $L^2(\Omega)$ 上的两个等价范数, \eqref{eq:ch1-negative-norm-pressure-estimate} 即得证.

为证明 $L^2(\Omega)$ 对范数 $[\sigma]$ 完备, 考虑关于该范数的 Cauchy 序列 $\sigma_m$. 则积分 $\int_\Omega\sigma_m\,dx$ 在 $\mathbb R$ 中构成 Cauchy 序列, 且导数 $\partial\sigma_m/\partial x_i$ 在 $H^{-1}(\Omega)$ 中构成 Cauchy 序列:
\begin{equation}
\int_\Omega\sigma_m\,dx\longrightarrow\lambda
\quad\text{当 }m\to\infty,
\label{eq:ch1-negative-norm-integral-limit}
\end{equation}
\begin{equation}
\frac{\partial\sigma_m}{\partial x_i}\longrightarrow\chi_i
\quad\text{当 }m\to\infty,
\quad\text{在 }H^{-1}(\Omega)\text{ 中},\quad1\leq i\leq n.
\label{eq:ch1-negative-norm-derivative-limit}
\end{equation}
显然 $(\operatorname{grad}\sigma_m,v)=0$, $\forall v\in\mathcal V$, 并且由 \eqref{eq:ch1-negative-norm-derivative-limit},
\[
\sum_{i=1}^{m}\langle\chi_i,v_i\rangle=0,
\qquad\forall v\in\mathcal V.
\]
由命题 \ref{prop:ch1-de-rham-gradient}, 存在某个分布 $\sigma$, 使
\[
\chi_i=\frac{\partial\sigma}{\partial x_i},
\qquad1\leq i\leq n.
\]
命题 \ref{prop:ch1-gradient-regularity} 表明 $\sigma\in L^2(\Omega)$. 可以选择 $\sigma$ 使
\[
\int_\Omega\sigma\,dx=\lambda,
\]
并且容易看出序列 $\sigma_m$ 在范数 $[\sigma]$ 下收敛到 $L^2(\Omega)$ 中的这个元素 $\sigma$.
\end{Proof}

\MBSourceStatementNumber{remark}{1}
\begin{Remark}
\label{rem:ch1-negative-norm-disconnected-domain}
若 $\Omega$ 不连通, 把 \eqref{eq:ch1-negative-norm-pressure-estimate} 中的 $\left\lvert\int_\Omega\sigma\,dx\right\rvert$ 替换为
\[
\sum_j\left\lvert\int_{\Omega_j}\sigma\,dx\right\rvert,
\]
则该不等式仍成立, 其中 $\Omega_j$ 是 $\Omega$ 的连通分支. 要把定理 \ref{thm:ch1-penalty-convergence} 推广到这种情形, 只需通过下列条件定义 $p$:
\begin{equation}
\int_{\Omega_j}p\,dx=0,
\qquad\forall\Omega_j.
\label{eq:ch1-penalty-component-pressure-zero-mean}
\end{equation}
\end{Remark}

\subsection{\texorpdfstring{$u_\varepsilon$}{u epsilon} 的渐近展开}
\label{subsec:ch1-penalty-asymptotic-expansion}

从现在起, 用 $u^0$ 和 $p^0$ 表示满足 \eqref{eq:ch1-penalty-pressure-zero-mean} 的 Stokes 问题 \eqref{eq:ch1-stokes-variational-problem}--\eqref{eq:ch1-stokes-weak-boundary} 的解(代替 $u$ 和 $p$). 将证明 $u_\varepsilon$ 有渐近展开
\begin{equation}
u_\varepsilon=u^0+\varepsilon u^1+\varepsilon^2u^2+\cdots
+\varepsilon^Nu^N+\cdots,
\label{eq:ch1-penalty-asymptotic-series}
\end{equation}
其中所有 $u^i$ 都属于空间 $\mathbf H_0^1(\Omega)$.

函数 $u^i$ 和一些辅助函数 $p^i$ 递归定义如下:
\begin{equation}
u^0,p^0\quad\text{已经给定};
\label{eq:ch1-penalty-base-pair}
\end{equation}
当 $u^{m-1},p^{m-1}$ 已知时($m\geq1$), 把 $u^m,p^m$ 定义为非齐次 Stokes 问题
\begin{equation}
u^m\in\mathbf H_0^1(\Omega),
\qquad p^m\in L^2(\Omega),
\label{eq:ch1-penalty-recursive-spaces}
\end{equation}
\begin{equation}
-\nu\Delta u^m+\operatorname{grad}p^m=0,
\label{eq:ch1-penalty-recursive-momentum}
\end{equation}
\begin{equation}
\operatorname{div}u^m=-p^{m-1},
\label{eq:ch1-penalty-recursive-divergence}
\end{equation}
\begin{equation}
\int_\Omega p^m(x)\,dx=0
\label{eq:ch1-penalty-recursive-pressure-zero-mean}
\end{equation}
的解. $u^m,p^m$ 的存在性与唯一性由定理 \ref{thm:ch1-nonhomogeneous-stokes} 得到. 条件 \eqref{eq:ch1-penalty-recursive-pressure-zero-mean} 有两个作用: 它保证 $p^m$ 完全唯一, 否则 $p^m$ 只在相差一个常数的意义下唯一; 它还保证层次 $m+1$ 所需的相容性条件:
\begin{equation}
\int_\Omega\operatorname{div}u^{m+1}\,dx
=\int_\Gamma u^{m+1}\cdot\nu\,d\Gamma
=0=\int_\Omega p^m\,dx.
\label{eq:ch1-penalty-recursive-compatibility}
\end{equation}

对 $N\geq1$, 记
\begin{equation}
u_\varepsilon^N=\sum_{m=0}^{N}\varepsilon^mu^m,
\label{eq:ch1-penalty-velocity-partial-sum}
\end{equation}
\begin{equation}
p_\varepsilon^N=\sum_{m=0}^{N}\varepsilon^mp^m.
\label{eq:ch1-penalty-pressure-partial-sum}
\end{equation}

\MBSourceStatementNumber{theorem}{2}
\begin{Theorem}
\label{thm:ch1-penalty-asymptotic-expansion}
设 $\Omega$ 是 $\mathbb R^n$ 中的有界 $C^2$ 类区域. 对每个 $m\geq1$, 存在由 \eqref{eq:ch1-penalty-recursive-spaces}--\eqref{eq:ch1-penalty-recursive-pressure-zero-mean} 唯一定义的函数 $u^m,p^m$.

对每个 $N\geq0$, 当 $\varepsilon\to0$ 时,
\begin{equation}
\frac{u_\varepsilon-u_\varepsilon^N}{\varepsilon^N}
\longrightarrow0
\quad\text{在 }\mathbf H_0^1(\Omega)\text{ 的范数中},
\label{eq:ch1-penalty-velocity-remainder}
\end{equation}
\begin{equation}
\frac{1}{\varepsilon^N}
\left(-\frac{\operatorname{div}u_\varepsilon}{\varepsilon}
-p_\varepsilon^N\right)
\longrightarrow0
\quad\text{在 }L^2(\Omega)\text{ 的范数中}.
\label{eq:ch1-penalty-pressure-remainder}
\end{equation}
\end{Theorem}

\begin{Proof}
存在性与唯一性已如前所述.

将 \eqref{eq:ch1-penalty-recursive-momentum} 乘以 $\varepsilon^m$, 再对 $m=1,\ldots,N$ 求和. 然后把所得方程与 $u^0,p^0$(以前记为 $u,p$)所满足的方程
\[
-\nu\Delta u^0+\operatorname{grad}p^0=f
\]
相加. 展开后得到
\[
-\nu\Delta u_\varepsilon^N
-\frac{1}{\varepsilon}\operatorname{grad}\operatorname{div}u_\varepsilon^N
=f-\varepsilon^N\operatorname{grad}p^N.
\]
与 \eqref{eq:ch1-penalty-velocity-convergence} 比较, 得
\begin{equation}
u_\varepsilon-u_\varepsilon^N\in\mathbf H_0^1(\Omega),
\qquad p^N\in L^2(\Omega),
\label{eq:ch1-penalty-remainder-spaces}
\end{equation}
\begin{equation}
-\nu\Delta(u_\varepsilon-u_\varepsilon^N)
-\frac{1}{\varepsilon}\operatorname{grad}\operatorname{div}
(u_\varepsilon-u_\varepsilon^N)
=\varepsilon^N\operatorname{grad}p^N.
\label{eq:ch1-penalty-remainder-equation}
\end{equation}
与 \eqref{eq:ch1-penalty-error-variational} 类似, \eqref{eq:ch1-penalty-remainder-equation} 等价于
\begin{equation}
\begin{aligned}
&\nu((u_\varepsilon-u_\varepsilon^N,v))
+\frac{1}{\varepsilon}
(\operatorname{div}(u_\varepsilon-u_\varepsilon^N),\operatorname{div}v)\\
&\hspace{14em}=-\varepsilon^N(p^N,\operatorname{div}v),
\qquad v\in\mathbf H_0^1(\Omega).
\end{aligned}
\label{eq:ch1-penalty-remainder-variational}
\end{equation}
在 \eqref{eq:ch1-penalty-remainder-variational} 中取 $v=u_\varepsilon-u_\varepsilon^N$, 得
\[
\begin{aligned}
&\nu\lVert u_\varepsilon-u_\varepsilon^N\rVert^2
+\frac{1}{\varepsilon}
\lvert\operatorname{div}(u_\varepsilon-u_\varepsilon^N)\rvert^2\\
&\quad=-\varepsilon^N(p^N,\operatorname{div}(u_\varepsilon-u_\varepsilon^N))\\
&\quad\leq\varepsilon^N\lvert p^N\rvert
\lvert\operatorname{div}(u_\varepsilon-u_\varepsilon^N)\rvert\\
&\quad\leq\frac{1}{2\varepsilon}
\lvert\operatorname{div}(u_\varepsilon-u_\varepsilon^N)\rvert^2
+\frac{\varepsilon^{2N+1}}{2}\lvert p^N\rvert^2,
\end{aligned}
\]
所以
\begin{equation}
\nu\lVert u_\varepsilon-u_\varepsilon^N\rVert^2
+\frac{1}{2\varepsilon}
\lvert\operatorname{div}(u_\varepsilon-u_\varepsilon^N)\rvert^2
\leq\frac{\varepsilon^{2N+1}}{2}\lvert p^N\rvert^2.
\label{eq:ch1-penalty-remainder-estimate}
\end{equation}
不等式 \eqref{eq:ch1-penalty-remainder-estimate} 显然蕴含 \eqref{eq:ch1-penalty-velocity-remainder}. 后者又蕴含 $\varepsilon^{-N}\Delta(u_\varepsilon-u_\varepsilon^N)\to0$ 在 $\mathbf H^{-1}(\Omega)$ 中成立, 从而 \eqref{eq:ch1-penalty-remainder-equation} 表明
\begin{equation}
-\frac{1}{\varepsilon^{N+1}}
\frac{\partial}{\partial x_i}
\operatorname{div}(u_\varepsilon-u_\varepsilon^N)
\longrightarrow\frac{\partial p^N}{\partial x_i}
\quad\text{在 }H^{-1}(\Omega)\text{ 中},
\quad1\leq i\leq n.
\label{eq:ch1-penalty-remainder-derivative-limit}
\end{equation}
但是
\[
\begin{aligned}
-\frac{1}{\varepsilon}\operatorname{grad}\operatorname{div}u_\varepsilon^N
&=-\frac{1}{\varepsilon}\sum_{m=1}^{N}\varepsilon^m
\operatorname{grad}\operatorname{div}u^m\\
&=\frac{1}{\varepsilon}\sum_{m=1}^{N}\varepsilon^m
\operatorname{grad}p^{m-1}
=\operatorname{grad}(p_\varepsilon^N-\varepsilon^Np^N),
\end{aligned}
\]
结合 \eqref{eq:ch1-penalty-remainder-derivative-limit}, 得
\[
\frac{1}{\varepsilon^N}\frac{\partial}{\partial x_i}
\left(-\frac{\operatorname{div}u_\varepsilon}{\varepsilon}
-p_\varepsilon^N\right)
\longrightarrow0
\quad\text{在 }H^{-1}(\Omega)\text{ 中},
\]
其中 $\varepsilon\to0$. 最后, \eqref{eq:ch1-penalty-pressure-remainder} 由 \eqref{eq:ch1-penalty-recursive-pressure-zero-mean}, \eqref{eq:ch1-penalty-pressure-partial-sum} 和引理 \ref{lem:ch1-negative-norm-pressure-estimate} 得到.
\end{Proof}

\MBSourceStatementNumber{remark}{2}
\begin{Remark}
\label{rem:ch1-penalty-asymptotic-disconnected-domain}
容易把注 \ref{rem:ch1-negative-norm-disconnected-domain} 用于定理 \ref{thm:ch1-penalty-asymptotic-expansion}: 若把条件 \eqref{eq:ch1-penalty-recursive-pressure-zero-mean} 替换为
\begin{equation}
\int_{\Omega_j}p^m\,dx=0
\label{eq:ch1-penalty-recursive-component-zero-mean}
\end{equation}
在 $\Omega$ 的每个连通分支 $\Omega_j$ 上成立, 则该定理对不连通集合 $\Omega$ 仍成立.
\end{Remark}

\MBSourceStatementNumber{remark}{3}
\begin{Remark}
\label{rem:ch1-penalty-discrete-asymptotic-version}
R.S. Falk [1] 研究了定理 \ref{thm:ch1-penalty-asymptotic-expansion} 所给结果的一个离散版本.
\end{Remark}

\subsection{数值算法}
\label{subsec:ch1-penalty-numerical-algorithms}

我们简要说明如何把第 \ref{sec:ch1-numerical-algorithms} 节所述算法推广到非齐次 Stokes 问题 \eqref{eq:ch1-penalty-recursive-spaces}--\eqref{eq:ch1-penalty-recursive-pressure-zero-mean} 的求解; 在这一阶段, 这是实际计算 $u_\varepsilon$ 的渐近展开 \eqref{eq:ch1-penalty-asymptotic-series} 的唯一困难.

这里只说明 Uzawa 算法的改造.

改变记号, 把问题 \eqref{eq:ch1-penalty-recursive-spaces}--\eqref{eq:ch1-penalty-recursive-pressure-zero-mean} 写成求 $v,p$, 使
\begin{equation}
v\in\mathbf H_0^1(\Omega),
\qquad p\in L^2(\Omega),
\label{eq:ch1-penalty-algorithm-spaces}
\end{equation}
\begin{equation}
-\nu\Delta v+\operatorname{grad}p=0,
\label{eq:ch1-penalty-algorithm-momentum}
\end{equation}
\begin{equation}
\operatorname{div}v=\phi,
\label{eq:ch1-penalty-algorithm-divergence}
\end{equation}
\begin{equation}
\int_\Omega p(x)\,dx=0,
\label{eq:ch1-penalty-algorithm-pressure-zero-mean}
\end{equation}
其中给定的 $\phi$ 满足
\begin{equation}
\int_\Omega\phi(x)\,dx=0.
\label{eq:ch1-penalty-algorithm-data-compatibility}
\end{equation}
$v,p$ 的存在性与唯一性已知.

以任意 $p^0$ 启动算法:
\begin{equation}
p^0\in L^2(\Omega),
\qquad\int_\Omega p^0(x)\,dx=0.
\label{eq:ch1-penalty-algorithm-initial-pressure}
\end{equation}
当 $p^m$ 已知时, 由下列方程定义 $v^{m+1}$($m\geq0$):
\begin{equation}
\left\{
\begin{aligned}
&v^{m+1}\in\mathbf H_0^1(\Omega),\\
&\nu((v^{m+1},w))-(p^m,\operatorname{div}w)=0,
\qquad\forall w\in\mathbf H_0^1(\Omega),
\end{aligned}
\right.
\label{eq:ch1-penalty-uzawa-velocity-step}
\end{equation}
\begin{equation}
\left\{
\begin{aligned}
&p^{m+1}\in L^2(\Omega),\\
&(p^{m+1}-p^m,\theta)
+\rho(\operatorname{div}v^{m+1}-\phi,\theta)=0,
\qquad\forall\theta\in L^2(\Omega).
\end{aligned}
\right.
\label{eq:ch1-penalty-uzawa-pressure-step}
\end{equation}
方程 \eqref{eq:ch1-penalty-uzawa-velocity-step} 是 $v^{m+1}$ 的 Dirichlet 问题:
\begin{equation}
v^{m+1}\in\mathbf H_0^1(\Omega),
\qquad-\nu\Delta v^{m+1}=-\operatorname{grad}p^m
\in\mathbf H^{-1}(\Omega),
\label{eq:ch1-penalty-uzawa-dirichlet-step}
\end{equation}
而 \eqref{eq:ch1-penalty-uzawa-pressure-step} 直接给出
\begin{equation}
p^{m+1}=p^m-\rho(\operatorname{div}v^{m+1}-\phi)
\in L^2(\Omega).
\label{eq:ch1-penalty-uzawa-pressure-update}
\end{equation}
注意到
\begin{equation}
\int_\Omega p^m\,dx=0,
\qquad\forall m\geq0.
\label{eq:ch1-penalty-uzawa-iterates-zero-mean}
\end{equation}

与定理 \ref{thm:ch1-uzawa-convergence} 完全相同(另见注 \ref{rem:ch1-uzawa-zero-mean-pressure}), 可以证明下述结果.

\MBSourceStatementNumber{theorem}{3}
\begin{Theorem}
\label{thm:ch1-penalty-uzawa-convergence}
若数 $\rho$ 满足
\begin{equation}
0<\rho<2\nu,
\label{eq:ch1-penalty-uzawa-step-condition}
\end{equation}
则当 $m\to\infty$ 时, $v^{m+1}$ 在 $\mathbf H_0^1(\Omega)$ 的范数中收敛到 $v$, 且 $p^{m+1}$ 在 $L^2(\Omega)$ 中弱收敛到 $p$.
\end{Theorem}
